% \iffalse meta-comment
%
% hit-thesis-toc.dtx 是目录与图表目录
%
% \fi
%
% \section{目录与图表目录}
%
%    \begin{macrocode}
%<*thesis-toc-lists>
%    \end{macrocode}
%
%
% 定义索引、目录格式
%
% \changes{v3.2a}{2026/08/08}{\cs{hit@start@toc} 与 \cs{hit@list@of} 迁到 expl3}
% \texttt{\#1} 是浮动体类型（\texttt{figure}、\texttt{table}、\texttt{equation}），
% 用来拼出 \cs{listfigurename}、\cs{ext@figure} 这些名字。
%    \begin{macrocode}
\cs_new_protected:Npn \hit@start@toc #1
  {
    \cs_set_eq:NN \oldnumberline \numberline
    \cs_set:Npn \numberline ##1
      { \oldnumberline { \use:c { #1 name } \skip_horizontal:n { .4em } ##1 } }
    \@starttoc { \use:c { ext@ #1 } }
    \cs_set_eq:NN \numberline \oldnumberline
  }
\cs_new_protected:Npn \hit@list@of #1#2
  {
    \hit@clear@page
    \hit@appendix@chapter * [ \use:c { list #1 name } ]
      { \use:c { list #1 name } } [ #2 ]
    \hit@start@toc {#1}
  }
%    \end{macrocode}
%
% 只有英文版用到，且需要clearpage
% \changes{v3.0.0}{2020/05/21}{只有英文版用到，且需要clearpage}
%    \begin{macrocode}


\cs_set:Npn \listoffigures   { \hit@list@of { figure } { \listfigureename } }
\cs_set_nopar:Npn \l@figure  { \addvspace { 6bp } \@dottedtocline {1} {0em} {4em} }
\cs_set:Npn \listoftables    { \hit@list@of { table } { \listtableename } }
\cs_set_eq:NN \l@table \l@figure
\cs_set:Npn \ext@equation    { loe }
\cs_set:Npn \equcaption #1
  { \addcontentsline { \ext@equation } { equation } { \protect \numberline {#1} } }
\cs_new:Npn \listofequations { \hit@list@of { equation } { \listequationename } }
\cs_set_eq:NN \l@equation \l@figure
%    \end{macrocode}
%
% \subsubsection{目录}
% \label{sec:toc}
% 本科文科生要求目录有四级。
% \changes{v3.2a}{2026/08/07}{目录相关代码从 bookcls 拆出至 book-toc 模块；拼接位置保留在原处以避免被后续包覆盖布局参数}
%    \begin{macrocode}
%</thesis-toc-lists>
%    \end{macrocode}
%
% ^^A ======================================================================
% ^^A Module thesis-toc: 目录格式（hit-thesis）
% ^^A ======================================================================
% \changes{v3.2a}{2026/08/11}{\cs{tocdepth} 改在 \cs{AtBeginDocument} 设，
%   \option{toc}/\option{depth} 写在导言区才来得及生效}
% \option{toc}/\option{depth} 数的是级数，\cs{tocdepth} 数的是比章低几级，两者差一。
% 这一句原先跟着类文件走，那时候导言区的 \cs{hitsetup} 还没执行，写了不起作用；
% 挪到 \cs{AtBeginDocument} 才赶得上，又仍旧早于第一条 \cs{addcontentsline}。
%    \begin{macrocode}
%<*thesis-toc>
\setcounter { secnumdepth } { 3 }
\AtBeginDocument
  { \setcounter { tocdepth } { \int_eval:n { \g__hit_toc_depth_int - 1 } } }
%    \end{macrocode}
%
% \changes{v3.2a}{2026/08/07}{book-toc 的加载期条件与 \cs{tableofcontents} 的
%   控制流改用 expl3。\cs{l@chapter} 里那串字体判断与 \cs{ctexset} 的取值保持
%   原样，它们是延后求值的排版内容，不是加载期逻辑}
% 工大论文目录中的潜规则：目录中的目录位置是空白。
% \begin{macro}{\tableofcontents}
%    \begin{macrocode}
\cs_set:Npn \tableofcontents
  {
    \bool_if:NTF \g__hit_openright_bool { \cleardoublepage } { \clearpage }
%    \end{macrocode}
%

% \changes{v3.1d}{2025/03/03}{如果是英文版，那只需要一个目录}
%    \begin{macrocode}
    \bool_if:NF \g__hit_en_bool
      {
        \phantomsection
        \markboth { \contentsname } { ccontent }
        \hit@chapter* { \hit@spread@by:nn { \hit@heading@spread } { \contentsname } }
        \pdfbookmark [ 0 ] { \contentsname } { ccontent }
        \normalsize \@starttoc { toc }
      }
    \__hit_tristate_if:NnTF \g__hit_toc_english_tl { \bool_if_p:N \g__hit_doctor_bool }
      { \tableofengcontents }
%    \end{macrocode}
%
% 有些硕士的范例中带有英文目录，这里添加手动开关
% \changes{v3.0.11}{2020/10/29}{有些硕士的范例中带有英文目录，这里添加手动开关}
%    \begin{macrocode}
      { }
%    \end{macrocode}
% \end{macro}
% 将英文目录潜入中文目录命令中，进一步简化thesis.tex结构
% \changes{v3.0.0}{2020/05/22}{将英文目录潜入中文目录命令中，进一步简化thesis.tex结构}
%    \begin{macrocode}
    \bool_if:NT \g__hit_en_bool
      {
        \printnoidxglossary
          [ sort=word , style=engstyle , title={\hit@nomenclature@heading} , nopostdot , nonumberlist ]
        \listoftables
        \listoffigures
      }
  }
%    \end{macrocode}
%
% 添加英文toc控制
% \changes{v3.0.0}{2020/05/21}{添加英文toc控制}
% 按照我工要求的目录格式。\cs{hit@toc@font} 在共享层定义，
% 排版时才读标志位，跟章节标题那边的 \cs{hit@title@font} 一样。
%    \begin{macrocode}
\cs_set:Npn \@pnumwidth { 4em }% 规定中的提前悬挂
\cs_set:Npn \@tocrmarg { \@pnumwidth }
\cs_set:Npn \@dotsep { 1 }
%    \end{macrocode}
%
% 此处临时更改一下对齐方式。\CTeX\ 似乎无法应对双语目录。
% todo:
% \changes{v1.0.04}{2017/09/17}{将leftskip设置参数至于外侧，以便后续添加可以适应标题长度的\cs{contentsline}方法}
%    \begin{macrocode}
\dim_set:Nn \@tempdima {4em}%
%    \end{macrocode}
%
% \changes{v3.2a}{2026/08/07}{\cs{@dottedtocline} 改成直接重定义，不再用
%   \pkg{etoolbox} 打补丁。补丁命令会把宏体 detokenize 后按当前 catcode 重扫，
%   在 expl3 里会吃掉宏体内的字面空格，只能退出 expl3 才敢用}
% 与内核原版比只有两处不同：条目文字外面套一层 \cs{hit@toc@font}，用
% \cs{use:c} 取；页码盒子从
% 定宽的 \cs{hb@xt@}\cs{@pnumwidth} 改成自然宽度的 \cs{hbox}，双语目录里页码
% 才不会被挤掉。
%    \begin{macrocode}
\cs_set:Npn \@dottedtocline #1#2#3#4#5
  {
    \int_compare:nNnF {#1} > { \c@tocdepth }
      {
        \vskip \z@ \@plus .2\p@
        {
          \leftskip #2\relax
          \rightskip \@tocrmarg
          \parfillskip - \rightskip
          \parindent #2\relax
          \@afterindenttrue
          \interlinepenalty \@M
          \leavevmode
          \@tempdima #3\relax
          \advance \leftskip \@tempdima
          \null \nobreak \hskip - \leftskip
          { \use:c { hit@toc@font } #4 } \nobreak
          \leaders \hbox { $ \m@th \mkern \@dotsep mu \hbox {.} \mkern \@dotsep mu $ }
          \hfill \nobreak
          \hbox { \hfil \normalfont \normalcolor #5 \kern -\p@ \kern \p@ }
          \par
        }
      }
  }
%    \end{macrocode}
%
% \changes{v2.0.08}{2017/06/15}{修改本科生论文目录格式（感谢QQ:嬴政　同学）}
% \changes{v3.0.01}{2020/06/06}{本科目录添加加粗}
% \changes{v3.0.0}{2020/05/21}{博后和英文版需要加粗}
% \changes{v3.2a}{2026/08/07}{目录里章条目的排版迁到 expl3，字体那串判断抽成
%   \cs{hit@toc@chapter@line@font}}
% 目录里章条目的字体。本科用 \cs{rmfamily}，其余取 \cs{hit@toc@font}
% （\option{toc}/\option{latin-sans} 关着的时候它展开成空）。之后按博后、
% 本科加粗、英文版依次叠加粗体。
%    \begin{macrocode}
\cs_new_protected:Npn \hit@toc@chapter@line@font
  {
    \bool_if:NTF \g__hit_bachelor_bool
      { \rmfamily } { \use:c { hit@toc@font } }
    \heiti
    \bool_if:NT \g__hit_postdoc_bool { \bfseries }
    \bool_lazy_and:nnT
      { \bool_if_p:N \g__hit_chapterbold_bool } { \bool_if_p:N \g__hit_bachelor_bool }
      { \bfseries }
    \bool_if:NT \g__hit_en_bool { \bfseries }
  }
%    \end{macrocode}
%
% 章条目本身。\cs{numberline} 在这里被调用，它要用到 \cs{@tempdima}。
% 标题与引导点之间那个 \texttt{\string~} 是原写法里换行带出来的空格，不能省。
%
% \changes{v3.2a}{2026/08/14}{\cs{@tempdima} 显式归零。原先没设，靠的是上一条
%   命令留在里面的残值，那个残值一变，章条目就逐字断行}
% \cs{@tempdima} 要自己归零。这一段读它两次——\cs{leftskip} 加一次，
% \cs{numberline} 用一次——却从没设过，用的是别处留在里面的值。
% \pkg{book} 的原版在这里写的是 \cs{setlength}\cs{@tempdima}\{1.5em\}，
% 迁写的时候漏了。漏了这些年没露馅是因为残值一直凑巧是零；新版面下版心一变，
% 残值成了版心高，\cs{leftskip} 跟着涨到一行只剩一个字宽，
% 章条目当场逐字断行。归零就是把一直在用的那个值写明白。
%    \begin{macrocode}
\RenewDocumentCommand \l@chapter { m m }
  {
    \int_compare:nNnT { \c@tocdepth } > { -1 }
      {
        \addpenalty { - \@highpenalty }
        \group_begin:
          \dim_zero:N \@tempdima
          \parindent \z@ \rightskip \@pnumwidth
          \parfillskip -\@pnumwidth
          \leavevmode
          \advance \leftskip \@tempdima
          \hskip -\leftskip
          { \hit@toc@chapter@line@font #1 } ~
          \leaders \hbox { $ \m@th \mkern \@dotsep mu \hbox { . } \mkern \@dotsep mu $ } \hfill
          \nobreak { \normalfont \normalcolor #2 } \par
          \penalty \@highpenalty
        \group_end:
      }
  }
%    \end{macrocode}
%
% 按工大标准, 缩小目录中各级标题之间的缩进，使它们相隔一个字符距离，也就是12pt。
% \changes{v3.1d}{2025/03/03}{根据 \issue[220] 来跟新目录中的编号与标题间距}
%    \begin{macrocode}
\cs_set_nopar:Npn \l@section       { \@dottedtocline {1} {1em} {2em} }
\cs_set_nopar:Npn \l@subsection    { \@dottedtocline {2} {2em} {3em} }
\cs_set_nopar:Npn \l@subsubsection { \@dottedtocline {3} {3\ccwd} {3.1em} }
%    \end{macrocode}
%
% 英文目录格式。
%    \begin{macrocode}
\cs_set:Npn~\@dotsep {~0.75~} % 定义英文目录的点间距
\dim_set:Nn \leftmargini {0pt}
\dim_set:Nn \leftmarginii {0pt}
\dim_set:Nn \leftmarginiii {0pt}
\dim_set:Nn \leftmarginiv {0pt}
\dim_set:Nn \leftmarginv {0pt}
\dim_set:Nn \leftmarginvi {0pt}

\cs_new:Npn \engcontentsname { \bfseries Contents }
\cs_new_protected:Npn \tableofengcontents
  {
~
%    \end{macrocode}
%
% 此处添加英文目录的章标题格式，默认细点
% \changes{v1.0.05}{2017/09/18}{矫正英文目录缩进与中文目录一致的bug}
%    \begin{macrocode}
 \cs_set_nopar:Npn \l@chapter {\@dottedtocline{0}{0em}{5em}}%控制英文目录： 细点\@dottedtocline  粗点\@dottedtoclinebold
 \@restonecolfalse
 \chapter*{\engcontentsname %chapter*上移一行，避免在toc中出现。
 \pdfbookmark[0]{Contents}{econtent}~
 \@mkboth{%
 \engcontentsname}{\engcontentsname}}~
%    \end{macrocode}
%
% 此处临时更改一下对齐方式。\CTeX\ 似乎无法应对双语目录。
% \changes{v1.0.04}{2017/09/17}{修正英文目录中换行时无法对齐的bug}
% 删除增加\cs{hangindent}的方法，其原因是\cs{numberline}多出一个空格
% \changes{v1.0.05}{2017/09/18}{彻底修正英文目录中换行时无法对齐的bug}
%    \begin{macrocode}
 \@starttoc{toe}%
 \if@restonecol\twocolumn\fi
  }
\cs_set:Npn~\@dotsep {~0.75~} % 定义英文目录的点间距
%    \end{macrocode}
%
% 目录中附录的章号格式。
%    \begin{macrocode}
\ctexset{%
  appendix/number=\hit@if@option@TF{bachelor}{\arabic{chapter}}{\Alph{chapter}},
}
%    \end{macrocode}
%
%    \begin{macrocode}
%</thesis-toc>
%    \end{macrocode}
%
% \Finale
